% ============================================================================
% SEC05.TEX
% Section 1.5: Tim Pengembangan Game
% ============================================================================

\section{Tim Pengembangan Game}

\begin{learningoutcomes}
\item Memahami struktur tim pengembangan game
\item Mengetahui peran dan tanggung jawab setiap anggota tim
\item Memahami pentingnya kolaborasi dalam tim
\item Mampu mengidentifikasi skill yang diperlukan untuk setiap peran
\end{learningoutcomes}

\subsection{Struktur Tim Pengembangan Game}

Tim pengembangan game dapat bervariasi tergantung ukuran project, dari tim kecil (indie) hingga tim besar (AAA studio).

\subsection{Peran Utama dalam Tim}

\subsubsection{1. Game Designer}

\textbf{Tanggung Jawab}:
\begin{itemize}
\item Merancang gameplay mechanics
\item Menyusun Game Design Document
\item Balancing gameplay
\item Prototyping mechanics
\end{itemize}

\textbf{Skills Required}:
\begin{itemize}
\item Kreativitas dan problem-solving
\item Pemahaman tentang player psychology
\item Kemampuan komunikasi dan dokumentasi
\item Basic programming knowledge (helpful)
\end{itemize}

\subsubsection{2. Programmer}

\textbf{Tanggung Jawab}:
\begin{itemize}
\item Implementasi gameplay mechanics
\item Engine programming
\item Tool development
\item Optimization dan debugging
\end{itemize}

\textbf{Specializations}:
\begin{itemize}
\item \textbf{Gameplay Programmer}: Core mechanics
\item \textbf{Engine Programmer}: Engine systems
\item \textbf{Graphics Programmer}: Rendering dan visual effects
\item \textbf{AI Programmer}: Artificial intelligence
\item \textbf{Network Programmer}: Multiplayer systems
\end{itemize}

\subsubsection{3. Artist}

\textbf{Tanggung Jawab}:
\begin{itemize}
\item Visual design dan art assets
\item Character design
\item Environment art
\item UI/UX design
\end{itemize}

\textbf{Specializations}:
\begin{itemize}
\item \textbf{Concept Artist}: Visual concepts dan ide
\item \textbf{2D Artist}: Sprites, textures, UI elements
\item \textbf{3D Artist}: Models, environments
\item \textbf{Animator}: Character dan object animation
\item \textbf{UI/UX Artist}: Interface design
\end{itemize}

\subsubsection{4. Audio Designer}

\textbf{Tanggung Jawab}:
\begin{itemize}
\item Music composition
\item Sound effects creation
\item Audio implementation
\item Voice direction (jika ada)
\end{itemize}

\subsubsection{5. Producer}

\textbf{Tanggung Jawab}:
\begin{itemize}
\item Project management
\item Timeline dan budget management
\item Team coordination
\item Stakeholder communication
\end{itemize}

\subsubsection{6. QA Tester}

\textbf{Tanggung Jawab}:
\begin{itemize}
\item Bug testing dan reporting
\item Gameplay testing
\item Compatibility testing
\item User experience testing
\end{itemize}

\subsubsection{7. Writer}

\textbf{Tanggung Jawab}:
\begin{itemize}
\item Narrative design
\item Dialog writing
\item Story development
\item World building
\end{itemize}

\subsection{Tim Kecil (Indie)}

Tim indie biasanya memiliki anggota yang memakai beberapa topi:
\begin{itemize}
\item Solo Developer: Semua peran
\item Small Team (2-5 orang): Generalist dengan spesialisasi minimal
\item Medium Team (5-15 orang): Mulai ada spesialisasi
\end{itemize}

\subsection{Kolaborasi dalam Tim}

\textbf{Pentingnya Komunikasi}:
\begin{itemize}
\item Regular meetings dan standups
\item Clear documentation
\item Version control untuk assets
\item Shared vision dan goals
\end{itemize}

\textbf{Tools untuk Kolaborasi}:
\begin{itemize}
\item Version Control: Git, Perforce, SVN
\item Project Management: Jira, Trello, Asana
\item Communication: Slack, Discord, Teams
\item Asset Sharing: Cloud storage, shared drives
\end{itemize}

\begin{catatan}
Dalam tim kecil, fleksibilitas dan kemampuan untuk memakai beberapa peran sangat berharga. Dalam tim besar, spesialisasi dan komunikasi yang efektif menjadi kunci sukses.
\end{catatan}

\subsection{Ringkasan Section}

\begin{itemize}
\item Tim pengembangan game terdiri dari berbagai peran spesialis
\item Peran utama: Designer, Programmer, Artist, Audio Designer, Producer, QA, Writer
\item Struktur tim bervariasi dari solo developer hingga tim besar
\item Kolaborasi efektif memerlukan komunikasi yang baik dan tools yang tepat
\item Dalam tim kecil, generalist skills sangat berharga
\end{itemize}
